\section{方阵乘法指数严格降低时的稳定排序}\label{sec:sorting}

张量网络有时须将一张表搬到使用处，有时须将对同一张表的所有查询集中到一起。搬运表涉及长记录，汇集查询则需要给许多短记录排序。我们将证明，只要指数严格低于 $9/4$，这两种操作的成本就低到足以完成第~\ref{sec:assembly}、\ref{sec:shaving} 节的最终组装。

固定第~\ref{sec:obstruction} 节的有限域 $k$，并令 $n=2^h$。第~\ref{sec:list} 节的带标记置换张量允许将每条记录的目的地作为输入的一部分。以 $j_i$ 表示宽 $w$ 比特的源记录 $x_i$ 的 $h$ 比特目的地址，以 $y_j$ 表示提议放在位置 $j$ 的记录。那么 $\mathrm{Tag}(h,w)$ 是下述事件的示性函数：各 $j_i$ 互异，且每个 $i$ 均满足 $y_{j_i}=x_i$。类似地，稳定排序张量表示以下事件：给定的候选输出已将输入记录按指定的键排序，且相同键的记录保持输入时的先后次序。两种示性函数的相等性检验都包括有效载荷的各个比特。

\begin{theorem}[条件稳定排序]\label{thm:sort} 设 $\omega_k<9/4$，并令 $n=2^h$。\begin{enumerate} \item 存在 $1<p<2$，使得 $\mathrm{Tag}(h,h)$ 在 $k$ 上有规模为 $O(nh^p)$ 的精确三价网络。对任意固定的 $c$，将宽度不超过 $ch$ 比特的记录作稳定排序也有相同的界。\item 对每个 $\zeta>0$，存在 $a>1$，使得对所有 $w\ge h^a$，$\mathrm{Tag}(h,w)$ 与宽 $w$ 比特记录的稳定排序在 $k$ 上都有规模为 $O_\zeta(nw h^\zeta)$ 的精确三价网络。\end{enumerate} \end{theorem}

\subsection{将查询排序以计算表值}

在元数为 $t$ 的站点，与之相接的各比特给出长度为 $2^t$ 的表中某一表项的地址。同时计算该站点的许多副本时，所选表项的乘积只取决于各个地址出现的次数。将查询排序即可求出这些次数，并将查表成本分摊到整个批次。

\begin{lemma}[批量查表]\label{lem:batch-table} 假设定理~\ref{thm:sort} 关于通常宽度的稳定排序界在 $k$ 上成立，指数为 $p>1$。设一个由二元线路与取值于 $K$ 的表组成的网络表示张量 $T$，其中 $K/k$ 为任意域扩张。若站点 $v$ 的元数为 $t_v\le u$，并且 $s=\sum_v(1+t_v)$，则 $T^{\otimes 2^u}$ 有一个精确的三价 $K$-网络，其规模为 \[
                  O\bigl(2^u s(1+u)^{p-1}\bigr).
\] 隐含常数与 $K$ 及各表的表项均无关。\end{lemma}

\begin{proof} 先考虑单张表 $f\colon\{0,1\}^t\to K$，其中 $t\ge1$，并将 $2^t$ 个查询地址 $x_1,\ldots,x_{2^t}$ 作为张量的外露腿。若地址 $z$ 出现 $c_z$ 次，所需乘积即为 $\prod_z f(z)^{c_z}$，其中未出现的地址对应空积。为每个 $z$ 创建一行字典记录，用 $t+1$ 个比特猜测其出现次数。在计数的第 $j$ 位（$0\le j\le t$）放置一个一元因子，其两个分量为 \[
                    (1,\ f(z)^{2^j}).
\] 该因子的位置已经确定了 $z$，因此两个分量都是网络的固定系数。

为每个查询建立单独的一行，并将其中的计数字段置零。将 $2^{t+1}$ 行先按地址、再按行的类型排序，使每一字典行位于对该地址的所有查询之后。计数器从零开始，每遇到一个查询便加一。遇到字典行时，要求计数器与该行携带的计数相符，再将计数器归零。即使一个地址没有查询，其计数也得到核验。对固定的外露地址，计数、排序后的行及计数器状态恰有一种合法取法。因此，在任意特征下，收缩的结果都等于 $\prod_i f(x_i)$，即使某个表项为零也成立。

用逐位进位加法实现计数器需 $O(2^t(1+t))$ 个门。将宽 $O(t+1)$ 比特的各行排序，成本为 $O(2^t(1+t)^p)$。排序网络在 $k$ 上实现一个零一张量，经标量扩张后便在 $K$ 上实现同一张量。

将每个元数为 $t$ 的站点的 $2^u$ 个副本分成 $2^{u-t}$ 批，每批 $2^t$ 个。元数为零的站点直接提供一个标量。在每个副本内依照原网络连接外露线路。所得网络收缩为 $T^{\otimes 2^u}$，其规模至多为下式的常数倍：\[
  2^u\sum_v(1+t_v)^p
       \ \le\ 2^u(1+u)^{p-1}\sum_v(1+t_v).
\] \end{proof}


\subsection{排序证明中的两种尺度}\label{sec:sort-road}

第二部分始终假设 $\omega_k<9/4$。我们先为宽记录安排路由。在一次跳转中，若只在较小的地址集内部置换记录，就可以递归处理。目标是把地址长度除以一个可任意指定的大常数，同时使每个求和项中传输的总宽度至多为输入宽度的常数倍，且此常数与缩短因子无关。重复这一归约即得定理~\ref{thm:sort} 对宽记录的估计。在此基础上，经过两个层次的基数排序即可得到通常宽度的估计。先通过短键在块内选出一个数位，再将该数位相同的记录整包传送。

第~\ref{sec:sort-action} 节用第一部分的关联不等式来比较目标置换的价格与较小地址集上的各次跳转的价格。这个比较首先对 CAT(0) 空间上的作用成立。只要局部价格数据足以实现这样的作用，就能得到相同的界。若数据未通过局部相容性检验，则可由谱比较得到一个降低成本的张量恒等式。第~\ref{sec:sort-route} 节结合这两种结论为有效载荷安排路由，并控制一个固定网络中的递归调用次数。

第二至四部分一律以二为底取对数来计量普通价格。在该单位下，距离 $d$ 即传输一个比特的价格。对固定的有限站点列表，由 \eqref{eq:prices-range} 和第~\ref{sec:prices-background} 节，尺度有界的价格数据（包括尺度为零的情形）构成紧集。在这里，给出直径上界就够了。沿二元线路的一棵生成树要求取值相等即可实现全对角张量，从而给出其价格的上界。

在固定比例处做有限次测地插值，所得站点间的距离连续依赖于原来的价格。扩充后的构形仍具有有界的尺度。沿任一收敛子网，插值等式在极限中仍成立，而此时引理~\ref{lem:geodesic} 可由原始数据确定各距离。

\section{从关联关系到置换作用}\label{sec:sort-action}

可将置换分解为每次只改变地址中一小段的若干次跳转。将其价格相加，得到与因子数成正比的界。为宽记录排序则要求界与因子数无关。我们先在度量空间中为位移求出这样的界，再证明只须对地址进行有限检验，就可控制有效载荷线路的平均价格。

\subsection{两个关联不等式}

固定 $k$ 的特征的一个幂 $q\ge100$。对一个二部关联图的配置，给两侧的每个顶点分别指定一个站点。均匀选取一对关联顶点，两端站点间距离的平方的均值称为关联能量 $e$。

\begin{lemma}\label{lem:sort-inc} 假设 $\omega_k<9/4$。在 $k$ 上任一完备 $\operatorname{CAT}(0)$ 谱空间中，考虑以下两种图之一的配置。\begin{enumerate} \item 射影空间 $\mathrm{PG}(\ell-1,q)$ 的点与直线，其中 $\ell\ge3$。\item $\mathbb F_Q^2$ 的点以及沿 $\mathbb F_q$ 上定义的 $q+1$ 个方向的仿射直线，其中 $Q=q^R$。\end{enumerate} 设 $o$ 为点侧各站点的重心。记 $A$、$B$ 分别为点侧和线侧各站点到 $o$ 的距离平方的均值，每一侧内部均匀加权。则 \[
                 e\ge (1-q^{-1/2})(A+B).
\] \end{lemma}

\begin{proof} 域 $\bar k$ 包含命题中出现的每个 $\mathbb F_Q$，并满足 $\omega_{\bar k}\le\omega_k$。要在这个域上应用定理~\ref{thm:incidence}，回顾标量扩张保持任一有限 $k$-张量配置的所有可能联合谱取值（第~\ref{sec:prices-background} 节）。特别地，我们的度量空间的每个有限子集都可等距实现在 $\bar k$ 上的某个完备谱空间中。对包含一个固定基点的有限子集，取这些实现的带基点超极限，所用超滤包含按包含关系排列的每个尾集。每个原有站点最终都会出现，并具有全部预定距离。由完备性和测地线的唯一性，所得等距像闭且凸。对超极限中的有界有限分布，重心方差不等式表明其重心等于各坐标重心的极限。若分布支撑于上述像内，由最近点投影可知其重心也在像内。因此只需对 $\bar k$ 上的配置证明这些不等式。

对仿射平面，将推论~\ref{cor:inc-link} 证明中的单位质量估计应用于构形 $\mathbb F_Q^2\subset\bar k^2$，保留 $q+1$ 个方向。距离平方比较式的交叉项由此至多为 $2\sqrt{AB/(q+1)}$，结论即得。

对射影空间，取一个从 $\bar k^\ell$ 到 $\bar k^3$ 的线性映射，使它保持由至多三个射影代表元组成的集合的秩。这样的映射构成一个非空开集。事实上，每个保秩条件都是非空开条件，而条件只有有限多个。在所得的射影平面中选择包含全部像点的仿射坐标片。二元组和三元组的秩不变，保证各像点互异，并且恰好保持原来所有共线的三元组。因此，这个仿射构形的割线代表原射影空间的直线。

每一点上有 $j$ 条直线，每一直线上有 $s=q+1$ 个点。对点站点记 $r_i=d(o,x_i)$，对线站点记 $r_L=d(o,y_L)$。令 \[
 K_L=\sup_{z\ne o}\Bigl(\sum_{i\in L}
       r_i\cos\angle_o(x_i,z)\Bigr)_+ ,
\] 其中 $r_i=0$ 时对应项取零。对单点空间结论显然成立。其余情形下，定理~\ref{thm:incidence} 给出 $\sum_L K_L^2\le (j-1)\sum_i r_i^2$。设 $M$ 为关联总数，由三角形比较和柯西--施瓦茨不等式，平均交叉项具有上界 \[
 \frac{2}{M}\sum_L r_L K_L
 \le 2\sqrt{\frac{j-1}{js}AB}
 \le q^{-1/2}(A+B).
\] 这里，$M$ 既等于点数的 $j$ 倍，也等于直线数的 $s$ 倍。从 $A+B$ 减去交叉项便得到结论。\end{proof}

\subsection{短地址集上的置换}

选定整数 $\ell\ge3$ 并预留 $m$ 个辅助坐标，其中 $m$ 将只依赖于 $q$。对 $n=2^h$，令 \[
 r=\left\lfloor \frac{\log_q n}{\ell+m}\right\rfloor,\qquad
 U=\mathbb F_q^r,\quad N=q^r,\quad
 c=\left\lceil n/N^{\ell+m}\right\rceil,\qquad
 \mathcal I=U^\ell\times[c]\times U^m .
\] 这里 $[c]=\{1,\ldots,c\}$。固定 $\mathcal I$ 上的一个次序，将原来的 $n$ 个地址与前 $n$ 个标签等同，并以增加不动点的方式扩充输入置换 $\pi$。补齐后的地址集满足 $n\le|\mathcal I|<2n$。

\emph{置换函数}以高度 $h$ 和一个合法的输入置换为参数，输出 $\mathcal I$ 上的一个置换。我们要求存在一个统一的确定性程序，使用 $nh^{O(1)}$ 次字操作，每个字宽 $O(h)$ 比特，通常的字操作均有多项式比特复杂度。这些函数在逐点复合下构成一个可数群 $G$。对显式列出的置换，可在线性时间内求其逆，且 $\pi$ 本身是 $G$ 的元素。

将 $U^\ell$ 与 $\mathbb F_q^\ell\otimes U$ 等同。固定维数至多为二的子空间 $V\le\mathbb F_q^\ell$，以 $K_V$ 表示保持投影 $\mathcal I\to(\mathbb F_q^\ell/V)\otimes U$ 的每一条纤维的子群。每一纤维均包含辅助标签的全部可能取值。其大小为 $cN^{m+\dim V}$，因而 $K_V$ 中的每次跳转均在大小至多为 $2^{th+O(1)}$ 的集合内实施置换，其中 \begin{equation}\label{eq:sort-hop-height}
                         t=\frac{m+2}{m+\ell}.
\end{equation} 加法常数来自 $c\le q^{m+\ell}$。增大 $\ell$ 就可缩短纤维内的地址。

子群 $K_V$ 生成 $G$。后续路由协议将用到可重排网络 \cite{benes} 的坐标分解的各个因子，我们在此给出算法。将 $[c]\times U^m$ 看作一个坐标。要消去大小为 $a$ 的坐标，以其余的输入、输出名称为两侧顶点构造二部多重图，每一项传输请求对应一条边。各顶点的度均为 $a$。用 $a$ 种颜色给边着色，每个顶点处每种颜色恰用一次。第一阶段，把所选坐标置换为相应边的颜色。最后阶段，将颜色置换成目的坐标。在这两个阶段之间，对每种颜色分别置换其余坐标。对固定的 $\ell,m$，迭代可得到有界个阶段。每个普通坐标阶段属于对应的坐标直线子群，每个辅助阶段属于 $K_0$。

为完整起见，说明这一着色可在确定性时间内完成，其耗时除对数因子外与边数成线性（另见 \cite{color}）。对度为偶数的正则二部多重图，沿欧拉回路交替着色，可将其平分为两个正则子图。对奇数度 $a$，先删去一个完美匹配。为求此匹配，设每一侧的顶点数为 $v$，取满足 $2^H>av$ 的最小 $H$。将各边的重数乘以 $\lfloor 2^H/a\rfloor$，并用一个任意虚拟完美匹配的若干副本将图补到度 $2^H$。虚拟边的总重数小于 $av<2^H$。反复将图平分为两个正则子图，保留虚拟边副本较少的一半，同数时以确定性规则选择。每次平分先平分偶数重数，再沿欧拉回路交替分配由奇数重数余下的副本，虚拟边与原有边分别标识。经 $H$ 次平分，保留的图度为一，且不含虚拟边。只须存储原边与虚拟匹配，每个重数只占一个字。删去匹配后剩余的偶数度可以减半，因此整个着色递归只有对数层数。

\subsection{位移的一致上界}

沿坐标分解将各位移相加所得的界依赖于 $\ell$。利用下面引理中的不动点，可以用同一个绝对常数控制任意跳转之积的位移。

\begin{lemma}\label{lem:sort-geometry} 设 $G$ 等距作用于完备 $\operatorname{CAT}(0)$ 空间 $Y$，且引理~\ref{lem:sort-inc} 的不等式对该空间的一切配置都成立。存在绝对常数 $C_0$，使得如果对所有 $g\in\bigcup_V K_V$ 均有 $d(x,gx)\le1$，那么 $G$ 固定一个距 $x$ 至多 $C_0/2$ 的点。特别地，$d(x,\pi x)\le C_0$。\end{lemma}

\begin{proof} 对每个旗 $0\subset L\subset P$ 取一个三角形，其中 $L$、$P$ 分别为 $\mathbb F_q^\ell$ 中的一维、二维向量子空间。采用 Garland 的链环能量论证及其度量不动点形式 \cite{garland,izekinayatani}。均匀选取一个旗，再均匀选取其中的一个顶点或一条边。一个配置给每个 $V$ 指定一个被 $K_V$ 固定的点 $y_V$。对每个 $V$，我们要另找一个这样的点 $z_V$，使得在均匀选取的经过 $V$ 的旗中，有 \begin{equation}\label{eq:sort-flag-link}
 \mathbb E\,d(z_V,y_W)^2\le \tfrac34\,
                           \mathbb E\,d(y_{W_1},y_{W_2})^2 .
\end{equation} 这里 $W_1,W_2$ 是与 $V$ 相对的两个顶点，$W$ 在二者中均匀选取。在 $V=0$ 处，各条对边对应于射影点 $L$ 与射影直线 $P$ 的关联。各 $y_L$ 的重心被 $K_0$ 固定。由引理~\ref{lem:sort-inc}，该重心到邻点的距离平方的均值至多是对边能量的 $1/[2(1-q^{-1/2})]$ 倍。在 $V=L$ 处，满足 $P\supset L$ 的 $y_P$ 的重心被 $K_L$ 固定。取 $y_0$ 为检验点，重心方差不等式给出系数 $1/2$。

\smallskip\noindent \emph{平面顶点。}当 $V=P$ 时，对边连接 $0$ 与一条直线 $L\subset P$。这里 $y_0$ 和各 $y_L$ 未必被 $K_P$ 固定。用这些站点在群作用下的像构造一个边能量相同的仿射关联图。该图点侧的重心在 $K_P$ 作用下将只移动很小的距离。

固定一组基和一个线性补空间，在模 $P$ 的每条纤维中沿 $P$ 选取坐标 $(x_1,x_2)\in U^2$。每一对坐标标识一个含 $cN^m$ 个标签的类。考虑各纤维中另一任意分划，要求各类的名称和基数与原来相同，且能对所有高度和输入在定义 $G$ 的时间界内统一算出。在每条纤维内排序并匹配相同的名称，即可用某个 $g\in K_P$ 将坐标分划变到这一分划。具体地，若原坐标名称为 $p$，搬运后的名称为 $p\circ g^{-1}$。给新分划指定站点 $g y_0$。任意另一个这样的 $g$ 与原来所取的元素只相差右乘 $K_0$ 的一个元素，因此此站点定义良好。将名称模 $L\otimes U$ 投影，可得一个更粗的分划，给它指定 $g y_L$。原投影分划的稳定子为 $K_L$。这两项指派均与 $K_P$ 的作用可交换。因此一个分划与其投影的站点距离恰为 $d(y_0,y_L)$。

令 $Q=q^R$，其中 $R\le m$ 将在下文选定。将 $R$ 个辅助 $U$-坐标作为 $\mathbb F_q$-线性映射 $u:\mathbb F_Q\to U$ 在 $\mathbb F_Q$ 的一组固定基上的取值。在一条纤维中均匀取一个标签，则 $u$ 均匀分布，且独立于 $(x_1,x_2)$。对每个 $(a,b)\in\mathbb F_Q^2$，构造以 $(x_1+u(a),x_2+u(b))$ 为名称的分划。该名称在 $U^2$ 上均匀分布。沿任一方向 $L$ 定义在 $\mathbb F_q$ 上的仿射直线，由 $u$ 的线性性，模 $L\otimes U$ 的名称作为 $(a,b)$ 的函数保持不变。这样就为所有点和所有保留的直线指定了站点。关联能量为 $e_P=\mathbb E_{L\subset P}d(y_0,y_L)^2$。对点侧重心 $b_0$，引理~\ref{lem:sort-inc} 给出 $(1-q^{-1/2})(A+B)\le e_P$。

\smallskip\noindent \emph{不动中心。}从点站点的分布中独立抽取 $d_*$ 个样本，构造逐次测地均值：第 $i$ 步向第 $i$ 个样本前进全程的 $1/i$。其最终值 $M_{d_*}$ 满足 \begin{equation}\label{eq:sort-mean}
                  \mathbb E\,d(M_{d_*},b_0)^2\le A/d_*.
\end{equation} 事实上，以上一步的均值为条件，距离平方的凸性以及 $b_0$ 处的方差不等式给出新期望平方误差的上界，即原误差的 $(1-1/i)^2$ 倍加上 $A/i^2$。

若对每一高度和每一输入，在每条纤维内，一组分划的联合名称均匀分布，就称该组分划独立。这是已经实现的分划组所具有的性质。选取一个充分大的绝对整数 $d_*$，并令 $m=R=4d_*$。从 $\mathbb F_Q$ 独立均匀抽取 $4d_*$ 个样本，它们在 $\mathbb F_q$ 上线性相关的概率至多为 $\sum_{j=0}^{R-1}q^{j-R}<1/(q-1)$。以至少 $1/2$ 的概率，抽到的 $2d_*$ 个仿射点的两个坐标构成 $4d_*$ 个线性无关的域元素。给定 $(x_1,x_2)$ 后，$u$ 在这些域元素上的取值仍联合均匀分布，所以这 $2d_*$ 个分划独立。

通过在每条纤维内匹配联合名称，群 $K_P$ 可在此长度的所有能高效指定的独立分划组上作传递作用。因此，前后两半各自均值间的距离对所有这些分划组都相同。由 \eqref{eq:sort-mean} 和三角不等式，该距离平方的无条件期望至多为 $4A/d_*$。因此，每一独立分划组中该距离的平方至多为 $8A/d_*$。同样的条件化论证给出一组 $d_*$ 个独立的仿射分划，其均值 $M$ 满足 $d(M,b_0)^2\le2A/d_*$。

固定 $g\in K_P$。给出 $M$ 的分划组与其在 $g$ 作用下的像可能相关。构造第三组分划，使之分别与前两组中的每一组联合独立。在每一纤维中，对每个标签用一条边连接其原名称与变换后的名称。所得二部图的两侧各有 $N^{2d_*}$ 个顶点，每个顶点的度为 $cN^{m+2-2d_*}$。该度数可被 $N^{2d_*}$ 整除。如上给边作正常着色，并将颜色分为 $N^{2d_*}$ 个等大集合。将每条边所属的颜色集合视为它在 $U^{2d_*}$ 中的第三个名称。在每个顶点处，各个第三名称出现的次数相等，两项所需的独立性因此均成立。上述着色也在定义 $G$ 的时间界内给出了第三组分划。

该组的均值到 $M$ 和 $gM$ 的距离均不超过 $\sqrt{8A/d_*}$。因此当 $d_*$ 充分大时，\[
 d(b_0,gb_0)\le 2\bigl(\sqrt{2A/d_*}+\sqrt{8A/d_*}\bigr)
                      \le 0.04\sqrt A
\] $b_0$ 的 $K_P$-轨道的外心 $z_P$ 被 $K_P$ 固定，且与 $b_0$ 相距至多 $0.04\sqrt A$。将一半质量分配在仿射关联图的每一侧，到 $z_P$ 的均方根距离至多为 $\sqrt{(A+B)/2}+0.04\sqrt A$。求此平均时，每个仿射点站点都由 $y_0$ 经群作用得到，每个方向为 $L$ 的仿射直线站点都由 $y_L$ 经群作用得到。因为 $z_P$ 为不动点，这些距离恰好分别是 $z_P$ 到 $y_0$ 和 $y_L$ 的距离。由 $(1/\sqrt2+0.04)^2/(1-q^{-1/2})<3/4$ 得到 $P$ 处的 \eqref{eq:sort-flag-link}。

\smallskip\noindent \emph{在旗上迭代。}将 $y_V$ 的初值取为 $K_Vx$ 的外心，它距 $x$ 至多为一。把每个 $y_V$ 替换为它与 $z_V$ 的中点，该中点仍被 $K_V$ 固定。记 $e$ 为一个旗中边长平方的均值。由测地线间距离的凸性，$V,W$ 处新端点间的距离平方至多为 $\bigl(d(y_V,z_W)^2+d(z_V,y_W)^2\bigr)/2$。因此对所有旗平均 \eqref{eq:sort-flag-link}，可得 $e_{\rm new}\le3e/4$。以随机选取的原邻点为中间点应用三角不等式，各顶点移动距离的均方根至多为 $(\sqrt e+\sqrt{3e/4})/2\le\sqrt e$。起初 $e\le4$。在有限加权度量积中，各步的移动距离可求和，极限中每条边的长度为零。旗图连通，故极限配置只有一个取值 $y_*$。所有 $K_V$ 因而 $G$ 都固定它。将各步移动距离相加，得 $d(x,y_*)\le1+2/(1-\sqrt{3/4})$。可取 $C_0=100$。\end{proof}

\subsection{置换轨道上的有限检验}

现在回到有限地址集。路由协议为每条有效载荷线路分别给定一个候选价格。我们将在半径有界的球上检验这些价格。为了比较重叠球上的价格实现，还须对其公共站点上的其他张量进行检验。

从 $\bigcup_V K_V$ 中取一个有限的对称元素集，其中含有 $\pi$ 的一个固定坐标分解的各个因子及其逆，称这些元素为\emph{核心生成元}。对每个输入，它们在 $\mathcal I$ 上定义一个无向图，标签到它在某个生成元下的像之间的跳转称为一条\emph{核心链路}。图的度不超过该集合的大小。每个 $i$ 都有一条依次沿上述因子到达 $\pi(i)$ 的路径，其长度以某个固定有理数 $a_{\mathrm p}$ 为界，与 $h$ 和输入无关。

四个检验参数分别是核心生成元、球半径、允许的最大站点数和最大局部维数。将每条有效载荷线路看作一种商品，给它单独编号。再固定坐标基；对于同一检验球内一列有序的互异站点，对其上每个具有允许格式的非零 $k$-张量也作编号。编号在全图统一赋予，使得两个球若含有相同有序站点上的同一张量，所见价格就相同。由于球的大小有界且域有限，共有 $O(|\mathcal I|)$ 个编号。

给每种商品指定非负价格。每条有效载荷线路的价格以 $a_{\mathrm p}$ 为上限。对其余每种格式，选一个统一的有理上限，足以在一个站点制备张量，并以二进制坐标沿球内的核心链路传输各条腿。对一个检验球，将球内所有商品的指定价格与其站点上的一组普通谱价格比较。误差取以下各量中的最大者：各商品的绝对偏差，以及两端均在球内的各核心链路的价格超出单位成本的正差。重复编号的商品分别与同一张量的谱价格比较。在所有普通价格（允许坍缩价格）上取得的最小误差称为检验的\emph{偏差}。此最小值可以达到。事实上，核心链路将球连通，因此一棵生成树即可控制任何误差有界的价格的尺度，而尺度有界的价格构成紧集。在单点上采用坍缩价格。如果此最小值至多为 $\rho$，就说该中心在容差 $\rho$ 下\emph{通过检验}。

\begin{lemma}[有限检测]\label{lem:sort-local-tests} 对每个 $\ell$，都存在一组有限检验参数和 $\rho>0$，使得对一切充分大的 $h$ 和任意合法的置换输入，以下结论成立。如果在容差 $\rho$ 下未通过检验的中心所占比例不超过 $\rho$，则指定的有效载荷价格满足 \[
        |\mathcal I|^{-1}\sum_{i\in\mathcal I} x_{i,\pi(i)}
                  \le C_0+1 .
\] 其中 $C_0$ 为引理~\ref{lem:sort-geometry} 的绝对常数。\end{lemma}

\begin{proof} 若结论不成立，选取一列满足 $h\to\infty$ 的反例。扩大核心生成元的有限对称集，使它们穷尽各 $K_V$ 的可数并，并令半径和格式上限趋于无穷。令容差和允许例外的比例趋于零，保持 $\pi$ 的原分解和有效载荷价格的上限不变。利用这些反例，我们将构造一个作用，使每次跳转的位移至多为一，而目标置换的位移至少为 $C_0+1$。

\smallskip\noindent \emph{标记站点。}对每个通过检验的中心 $i$，取一组达到检验最小误差的价格，并取其完备 $\operatorname{CAT}(0)$ 谱实现。对其余中心使用单点空间。对 $g\in G$，在中心 $i$ 的局部实现中标记站点 $g^{-1}i$。由于 $g^{-1}$ 可写成子群元素的一个固定的字，这些跳转最终都会纳入检验，且球中包含从 $i$ 到 $g^{-1}i$ 的整条路径。因此对每个固定的 $g$，从中心到这个标记的距离一致有界。如果标记不可用，就用中心自身。

在序列的每个指标处，取所有中心所对应空间的度量积，采用均方根距离，以各中心组成的向量为基点。对这些积取带基点的超极限，令 $Y$ 为所得标记的闭凸包。由中点不等式可知，积和超极限都是完备 $\operatorname{CAT}(0)$ 空间。引理~\ref{lem:sort-inc} 的两个关联界在 $Y$ 中均成立。积空间的重心可逐坐标求取，故可将各坐标上的平方不等式取平均。在超极限中，方差不等式表明重心的极限就是相应有界有限配置的重心。最后，对支撑于某个闭凸子集的分布，到该子集的投影保持其重心。

\smallskip\noindent \emph{标记的平移。}我们断言左乘 $g\mapsto ug$ 可延拓为 $Y$ 的等距变换。将标记作有限次测地插值所得的表达式给出 $Y$ 的一个稠密子集中的点。固定两个这样的表达式。在中心 $i$，平移后的表达式所对应的地址标签与原表达式在中心 $u^{-1}i$ 的地址标签相同，因为 $(ug)^{-1}i=g^{-1}u^{-1}i$。

当指标足够大时，两个球都包含这些标签。由于 $i\mapsto u^{-1}i$ 保持均匀测度，除去比例趋于零的例外，两个中心均通过检验。对互异标签上的每个固定张量，一旦格式上限将其纳入，两处实现的取值之差便至多为容差的两倍，因为二者使用同一商品编号。因此两个局部实现中这两个表达式间的距离之差，在上述中心对上一致趋于零。事实上，若有违反此结论的序列，可选取一个子列，使标签间的重合模式固定。有限多个因子字给出两个局部实现中标记列表的直径上界，故生成树估计控制了其价格的尺度。由紧性，取该列表上的极限价格。由于格式逐步穷尽所有张量，两套价格在每个张量上相等。由引理~\ref{lem:geodesic}，测地表达式间的距离连续，所以极限距离也相等，矛盾。

即使在例外中心对上，所有表达式间距离也是有界的。在通过检验的局部实现中，固定的字长控制所有标记的位置，而未通过的局部实现已经坍缩为一点。因此例外中心对对平方均值的贡献趋于零。按 $i\mapsto u^{-1}i$ 重新编号，就证明了极限中所断言的距离相等。由此，表达式上的映射可延拓到 $Y$。其逆由 $u^{-1}$ 给出，群运算律先在表达式上成立，再由连续性在 $Y$ 上成立。

令 $x$ 为单位元所对应的标记。对每个 $g\in K_V$，相应跳转及其逆最终均在检验集内，因此 $d(x,gx)\le1$。在中心 $i$，标记 $\pi^{-1}x$ 对应于地址 $\pi(i)$。在通过检验的局部实现中，此标记到中心的距离与有效载荷线路 $i\to\pi(i)$ 的指定价格相差不超过容差。在未通过的局部实现中，距离为零，价格至多为 $a_{\mathrm p}$。因此 $d(x,\pi^{-1}x)$ 即为有效载荷价格均方根的超极限。该超极限至少等于反例算术平均值的极限下界 $C_0+1$。由于 $d(x,\pi^{-1}x)=d(x,\pi x)$，这与引理~\ref{lem:sort-geometry} 矛盾。\end{proof}

\section{有界宽度的置换路由}\label{sec:sort-route} \subsection{从价格检验失败到路由恒等式}

固定引理~\ref{lem:sort-local-tests} 中的核心生成元及检验参数。记 $M=|\mathcal I|$，于是 $n\le M<2n$。沿指定路径从 $i$ 向 $\pi(i)$ 发送一个比特，就要为路径上的每次跳转付费。我们要寻求网络的一个精确和，使每个有效载荷比特的平均传输成本有一个绝对上界，即使参数 $\ell$ 增大也如此。

\begin{proposition}\label{prop:sort-protocol} 假设 $\omega_k<9/4$，固定 $H\ge1$。对一切充分大的 $h$、每一组合法标记和所有 $h\le w\le h^H$，相应的 $w$ 比特有效载荷置换张量都恰好等于 $k$ 上若干个网络的和，这些网络具有以下性质。\begin{enumerate} \item 对每个核心生成元，每个标签都有有界个固定槽位。启用的槽位沿生成元或其逆向一个邻点传输 $O(w)$ 个比特。两端属于同一个预先指定的组，该组至多含 $2^{th+O(1)}$ 个标签。任一求和项中所传输的宽度之和至多为 $Cnw$，其中 $C$ 与 $\ell$ 和 $t$ 无关。\item 在每个标签上，由规模为 $O(w(1+\log w)^{O(1)})$ 的固定三价网络实现局部映射。其配置由一个长度有界的列表组成，列中记录选单编号、宽度与偏移量。\item 一个求和项可用 $O(n\log(2w))$ 个比特指定。给定这些比特与标记，一个确定性控制器可用 $nh^{O(1)}$ 次字操作算出局部配置与匹配。各指数及槽位容量均与 $H$ 无关。运行时间可含有依赖于 $H$ 的乘法常数。\end{enumerate} \end{proposition}

\begin{proof} 我们使用 \cite[Section~\ext{G}{sec:protocol}]{oaG} 的局部替换论证。一种商品的一个\emph{请求}表示相应张量有待接到指定的张量腿上。其价格上限 $a_i$ 控制直接供给的成本，即在某一站点制备该张量，并沿指定的核心路径传输各条腿所需的成本。以二进制编码每条腿，将未用坐标补零。由多重线性性，在收缩的任一步都可用一个张量恒等式替换若干请求。

\smallskip\noindent \emph{球上的恒等式。}设价格数组 $x\in\prod_i[0,a_i]$ 在某球上的偏差为正。对球内每种满足 $x_i>0$ 的商品各请求一份。由于零数组能够通过检验，这些商品的张量积 $D$ 的指定总价格满足 $L=\sum_{i:x_i>0}x_i>0$。可用的供给 $T_j$ 包括按指定价格供应的球内各商品，以及价格为一的球内各二元核心链路。在单点处，所有需求都可以就地满足。在较大的球上，则存在严格高于各供给价格的价格 $l_j$，以及数 $0<L_-<L$，使得每一归一化谱对数价格 $P$ 均满足 \[
                 P(D)<L_-\max_j\frac{P(T_j)}{l_j}.
\] 事实上，二元核心链路的一棵生成树即可提供全二元对角张量，其归一化对数价格为一。因此最大比值不小于沿树的 $l_j$ 之和的倒数。若对从上方趋于原供给价格的扰动，以 $L$ 为系数的严格不等式不成立，则将每一个违反不等式的 $P$ 除以其最大比值以重新定标。生成树的界使这些尺度保持有界。于是一个极限价格满足所有供给上界，并且赋予 $D$ 至少为 $L$ 的价格。价格在乘积上的可加性迫使每种正价格商品处均取等号，其余商品的价格均为零。这个极限就会通过该球上的检验。固定一个使严格不等式成立的扰动，紧性保证可以把 $L$ 降至 $L_-$，同时保留一致的严格余量。

选取整数 $b_0$ 并构造直和 \[
               U=\bigoplus_j T_j^{\otimes\lfloor b_0L_-/l_j\rfloor}.
\] 取整只会在对数价格上造成有界损失。若 $b_0$ 足够大，则每个归一化谱价格在 $U$ 上的值都高于在 $D^{\otimes b_0}$ 上的值。由第~\ref{sec:prices-background} 节的严格谱比较，对每个 $r\ge1$，从至多 $K(b_0)$ 份 $U^{\otimes r}$ 的直和经精确限制即可得到 $D^{\otimes b_0r}$。展开各直和并应用限制映射。所得普通和满足每种需求的 $b=b_0r$ 个副本。每一项所用供给的总价格至多为 $bL_-$。该项的编号指定了一个直和副本及一个由 $r$ 次供给选择组成的有序字。若供给类型数为 $s$，这至多需要 $\lceil\log_2(K(b_0)s^r)\rceil$ 个比特。先增大 $b_0$ 再增大 $r$，可使该长度与 $b$ 之比任意小。

将此恒等式重复 $J$ 次。仅用 $J$ 个项编号的\emph{频数直方图}作外层选择器。由于字母表固定，该选择器只需 $O(\log(2J))$ 个比特。在相应网络内部，一个根节点对有序字求和，沿核心树传送各个字，并强制其频数直方图与指定值相符。直方图确定每一种供给的副本数，有序字则把这些副本指派给限制映射的各次应用。每个有序字都恰好贡献一次，其系数保持为 $k$ 中的原系数。将项编号的编码率选得足够低，使得沿树传输的成本至多为 $Jb(L-L_-)/2$。这样，即使计入替换本身所用的控制比特，每个需求块仍节省一个正量的成本。

在原数组的某个邻域内，这些需求的价格保持为正，替换也始终带来严格节省。由于域、半径、度数及格式均固定，球的类型只有有限种。偏差至少为 $\rho/2$ 的价格数组构成紧集，因此只需有限张替换选单。在该集合上，每单位需求所节省的成本以及选中需求的价格都有一致的正下界。

\smallskip\noindent \emph{控制总宽度。}考察展开式中的一项。令 $wz_i$ 表示该项对商品 $i$ 的待满足请求的整数存量，令 $S$ 表示已设置的核心线路宽度除以 $Mw$ 的商。将该项的账面值记作 $S+F(z)$，其中 \[
       F(z)=\frac1M\max_{0\le x\le a}
                   \sum_i\bigl(z_ix_i-\tfrac12\mu x_i^2\bigr).
\] 坐标数为 $O(M)$，各价格上限一致有界。选择正有理数 $\mu$，使得 $\mu\sum_i a_i^2/(2M)\le1$。使目标取最大值的价格 $x_i=\min(a_i,z_i/\mu)$ 满足 $z_i\ge\mu x_i$。由它的 Lipschitz 界，当 $z+v\ge0$ 时有 \[
       F(z+v)-F(z)
       \le \frac1M\sum_i x_iv_i+\frac1{2\mu M}\sum_i v_i^2 .
\] 固定整数 $B>\max_i\mu a_i$。存量一旦超过 $wB$，便直接供给超出的部分。在区间 $z_i\ge B$ 上边际价格等于 $a_i$，因此该操作不增加账面值。

若 $M$ 条有效载荷线路的最优价格的均值超过 $C_0+2$，则由引理~\ref{lem:sort-local-tests}，有多于 $\rho M$ 个球的偏差大于 $\rho$。用贪心装填法可选出 $\Omega(M)$ 个互不相交的这样的球，并在检验中留出固定余量。事实上，偏差关于价格数组是 $1$-Lipschitz 的。取足够细的有理网格，对与偏差至少为 $\rho$ 的集合相交的每一网格单元，都可指定一个适用于整个单元的有效选单条目。利用这些有限表便足以选取并装填合格的球。

固定一个小的正有理数 $\alpha>0$，在每个选中的球上取 $J=\lfloor\alpha w/b\rfloor$ 并施行相应替换。由 $z_i\ge\mu x_i$ 及所选需求价格的正下界，存量充足。把每条新的核心链路计入 $S$，其他供给则登记为新请求。在账面值的线性估计中，每个球每消去一个需求块，即使计入传送有序字的成本，也能节省固定的正量。所以当 $w$ 足够大时，对某个 $c>0$，这些球至少可在归一化账面值上节省 $c\alpha$。选单有限且各球不交，故 $\max_i|v_i|=O(\alpha)$，从而二次误差为 $O(\alpha^2)$。将 $\alpha$ 固定得足够小，使得每轮并行替换在产生的每一项中均令账面值至少降低一个固定的正量。一轮结束后，直接供给超过存量上限的部分。

从零开始，在固定的细网格上，以等额增量引入每条有效载荷线路的全部 $w$ 个副本，取整时使总数恰为 $w$。每次增加前反复执行并行轮，直到各有效载荷线路的平均价格至多为 $C_0+2$。账面值非负。各增量的线性总成本至多为 $C_0+2$，若网格足够细且 $w$ 足够大，则它们的二次成本之和至多为一。每次增加后直接供给超出上限的存量。故每一项只经历有界轮数，最终账面值至多为 $C_0+3$。直接满足所有剩余请求。在 $F$ 的定义中取 $x=a$，得 \[
             \frac1M\sum_i a_i z_i\le F(z)+1.
\] 因此总宽度至多为 $(C_0+4)Mw\le 2(C_0+4)nw$。最后这个系数是绝对常数。

\smallskip\noindent \emph{实现各求和项。}每一轮中，每个标签至多属于一个选中的球。经由某条给定边的每一项直接供给都起始于距该边有界的范围内。由于度数有界、每个邻域中的商品类型数有界且存量有上限，每个生成元只需有界个宽 $O(w)$ 的槽位。生成元确定包含该槽位两端的组。

在同一商品涉及的各站点上，将其请求按共同的次序排列，从队首移出已满足的请求，在队尾追加新请求。对每次替换，直方图确定各列表的长度。在其网络内部，用有序字的前缀计数为各次重复应用分配供给位置。对宽 $O(1+\log w)$ 比特的地址实施比较排序，每个站点使用 $O(w(1+\log w)^{O(1)})$ 个三价门便可实现这些置换和部分匹配。以确定性方式补全各部分匹配，并将空闲坐标置零。每个单独的块映射在 $k$ 上都具有固定规模，可用常数规模的三价网络实现。只需有界个局部字段来选定映射、长度和偏移量。

控制器枚举有界邻域，根据存量计算价格，用有限表选出互不相交的合格球，并指派各个槽位。在有界轮数内共有 $O(n)$ 个直方图，每个直方图都有唯一的 $O(\log(2w))$ 比特编码。要求编码合法，而且选择器中未使用的位置一律置零。只有局部计数和编号进入控制器，而有序字在网络内部求和。对 $h$ 比特的标记和 $O(\log(2w))$ 比特的计数进行枚举和存量更新，共需 $nh^{O(1)}$ 次字操作。这就证明了三项估计。\end{proof}

\subsection{从局部传输到稳定排序}\label{subsec:sort-recursion}

记 $S(h,w)$ 为 $\mathrm{Tag}(h,w)$ 在 $k$ 上的最小三价网络规模，并令 $n=2^h$。命题~\ref{prop:sort-protocol} 控制了每个求和项所用的传输量。要推导 $S$ 的递推式，须将所有求和项中的调用安排在同一个网络中。只要来自不同项的数据可以用较低成本送达同一调用，它们就能共用这次调用。

\paragraph{共用一组固定的调用。}固定维数 $\ell$、地址缩短因子 $t=(m+2)/(m+\ell)$，以及命题~\ref{prop:sort-protocol} 的协议。在每个项中，该协议至多有 $Ln$ 条链路，总宽度至多为 $Cnw$。这两个界中的常数 $C$ 与 $\ell$ 无关。选取满足 $L\le C/Y$ 的 $0<Y\le1$，将每条链路切分为宽 $w'=\lceil Yw\rceil$ 比特的小段。末段的填充比特在两端均固定为零。每一项的非空小段数至多为 \[
             Cnw/w'+Ln\ \le\ m_*n,\qquad m_*=2C/Y .
\] 对此固定协议，共有 $O(n)$ 个候选槽位，依照可以在其内部通信的短地址集分组。若为每组所处理的有界个链路端点或小段复制该组，其规模只会改变一个常数倍。

使用 \cite[Lemma~\ext{G}{lem:slack-compaction}]{oaG} 的分配图。这种图把至多 $cK$ 个候选连接到 $K$ 个槽位。任一大小至多为 $K/2$ 的活动候选子集都可单射地分配到相邻的槽位。各度数仅依赖于 $c$，当图作固定辅助信息给定时，此分配需要 $O_c(K)$ 次确定性字操作。

选择一个二的幂次容量 $U=2^{th+O(1)}$，至少为最大候选组大小的两倍。用贪心法将整组装入容量为 $U$ 的列表；若下一组会使当前列表超过容量，就封闭该列表。每个已封闭的列表至少半满，故共有 $O(n/U+1)$ 个列表。固定一个较小的二进分数 $K/U$。在一个列表内，将活动小段按次序划为每批 $K/2$ 段，最后一批可较短。对每一非空批预留 $K$ 个槽位，则总预留数至多为 \[
             2m_*n+O((n/U+1)K)\ \le\ 3m_*n .
\] 这里 $U=o(n)$，且在装填估计的常数确定后再选 $K/U$。

对每一列表，设置足够多的候选批以容纳其中的每一小段。用一个分配图把一批的源端放入 $K$ 个槽位，再用另一个分配图处理目的端。第二步的分配对象是整批数据，目的地为共用的 $\lceil6m_*n/K\rceil$ 个物理模块，每个模块有 $K$ 个输入位置和同样多的输出位置。至多半数模块被激活，候选批数至多为模块数的常数倍。因此每个活动批的两张列表可同时移到为它分配的模块。计入空闲模块，对一切充分大的 $h$，总共设置的位置不超过 $7m_*n$ 个。

每个模块使用 $\mathrm{Tag}(h',w')$，其中 $h'=\log_2 K=th+O(1)$。顺次填补未匹配的位置，将未用比特固定为零，给空闲模块配置恒等置换的标记与全零有效载荷。一旦给定求和项及其标记，所有分配和填充均以确定性方式完成。由于分配图的度数有界，开关的总成本为 $O(nw)$，其中包括一次搬运两张各宽 $Kw'$ 比特的列表的开关。

其余控制只使用标记、选择器和槽位描述。可通过猜测一份带时间戳的执行记录来核验控制：逐条验证指令，将全部内存操作依地址、时间和同一指令内的访问次序作比较排序，以核验每次读取都与上次写入相符。固定的初始内存与填充使每一合法选择器恰有一条通过检验的执行记录。另用常规比较网络检查原标记是否互异，并拒绝不合法的置换。对每个固定的 $H$ 以及 $h\le w\le h^H$，较短地址调用以外的规模为 \begin{equation}
  D_{\rm control}(h,w)=
  O_H\bigl(nw(\log(2h))^{O(1)}+n h^J\bigr),\qquad J>2 .
  \label{eq:sort-control}
\end{equation} 第一项计入局部映射及其关于 $\log(2w)$ 的多项式开销。第二项计入 $nh^{O(1)}$ 次确定性字操作、其比特实现和内存核验。各指数均可取为与 $H$ 无关。由此得到递推式 \[
 S(h,w)\ \le\ D_{\rm control}(h,w)
       +\frac{7m_*n}{2^{h'}}\,S(h',w').
\]

\paragraph{宽有效载荷。}只需在 $0<\zeta<1$ 时证明估计。取 $\ell$ 充分大，使 $14C t^\zeta<1/4$，随后固定所得的 $Y$ 和 $J$。地址从 $h$ 缩短到 $h'$，可将允许的开销 $h^\zeta$ 乘以一个趋于 $t^\zeta$ 的因子。为处理宽度接近常规排序适用阈值的调用，我们让开销指数随宽度逐渐下降。

选取光滑非增函数 $u:[1,\infty)\to[\zeta,2]$，使其在 $[1,J]$ 上等于 $2$，并最终恒等于 $\zeta$。取足够长的过渡区间，使得 $b=-u'$ 处处满足 $0\le b<1$ 和 $Y^{-b}<3/2$。将 $H$ 固定在常值尾段的内部。对 $h\ge2$ 和 $h\le w\le h^H$，我们断言 \begin{equation}
       S(h,w)\ \le\ D n w h^{u(a)},\qquad a=\log_h w .
       \label{eq:sort-width}
\end{equation} 若 $a\le J$，则 Batcher 的双调排序网络~\cite{batcher} 可使用 $O(nh^2)$ 个比特成本为 $O(h+w)$ 的比较器，把带标记的源记录按目的地址排列。将其标记与规定位置比较，再将其有效载荷与候选输出比较，即得 \eqref{eq:sort-width}。

当 $a>J$ 时，对 $h$ 归纳。高度充分大时，$h'<h$ 且 $w'\ge h'$。记 $a'=\log_{h'}w'$，在 $J\le a\le H$ 上一致地有 \[
 a'-a=\frac{\log Y-a\log t+o(1)}{\log h},
 \qquad
 \frac{w'(h')^{u(a')}}{w h^{u(a)}}
   =(1+o(1))\,t^{u(a)+a b(a)}Y^{1-b(a)}.
\] 第二式是紧区间上的泰勒展开。若 $w'>(h')^H$，将有效载荷均分成 $\lceil w'/\lfloor(h')^H\rfloor\rceil$ 片，每片使用同样的标记。片数有界，各片宽度的总和等于 $w'$，宽度指数均为 $H+o(1)$，位于 $u=\zeta$ 的区域内。为每个槽位的每个切片复制标记的成本为 $O(h')$，已包括在 \eqref{eq:sort-control} 内。因此每个槽位仍满足同一估计。

乘以已设置位置数与原位置数之比 $7m_*=14C/Y$ 后，递归调用所占的成本至多是所断言上界的 $(1+o(1))14C t^\zeta Y^{-b(a)}<1/2$。由于 $a\ge J$ 且 $u\ge\zeta$，额外成本 \eqref{eq:sort-control} 一致为 $o(nwh^{u(a)})$。增大 $D$ 以涵盖有界的高度，即可完成归纳。对 $w>h^H$，将有效载荷分成尽可能均匀的切片，最大宽度为 $\lfloor h^H\rfloor$。当 $h$ 充分大时，各片的宽度指数均为 $H+o(1)$，复制标记的成本至多为总数据量的常数倍。对有界高度，双调网络的成本为 $O_h(w)$。这就证明，对某个 $A>1$，有 \begin{equation}
         S(h,w)=O_\zeta(2^h w h^\zeta)\qquad(w\ge h^A).
         \label{eq:sort-all-wide}
\end{equation}

\paragraph{从带标记置换到稳定排序。}对 $2^s$ 条宽 $v$ 比特的记录，稳定排序使用一次 $\mathrm{Tag}(s,v)$ 和额外的 $O(2^s(s+v))$ 个门。对每个候选输出，猜测它在原输入列表中的一个目的位置。用带标记置换核验每个输出都被送回一条与之相等的输入记录。再沿输出列表检查二元组（键，猜测的输入位置）按字典序严格递增。这样，具有同一键的所有记录必定恰好组成相应的输入子序列。即使记录完全相同，这一条件也唯一确定排列次序和各个猜测的下标。因此稳定排序关系继承宽记录情况下的界。

\paragraph{打包通常宽度的记录。}在基数排序的每一趟中两次使用宽记录的界。先在 \eqref{eq:sort-all-wide} 中取 $\zeta=1/2$，相应阈值为 $a>1$；再在同一式中取 $\zeta=1/(2a)$，相应阈值为 $B>1$。起初共有 $O(n)$ 条记录，每条记录的键和有效载荷均为 $O(h)$ 比特。选择宽 $s=\lfloor h^{1/a}\rfloor$ 比特的基数位。从最低有效位起实施 $O(h/s)$ 趟稳定排序，最后一个数位不足之处补零。两种排序的地址长度将分别约为 $2s$ 和 $h$。

对其中一趟，先填充列表，使每块都有 $2^{2s}$ 个槽位。在每一块内，按虚拟标志和当前数位排序，将虚拟记录按原次序排到最后。选取二的幂 $P\in[h^{B-1},2h^{B-1})$。对每一个由同一数位的真实记录构成的连续段，在末尾插入虚拟记录，使段长成为 $P$ 的整倍数。现在每连续 $P$ 个槽位组成的一包中，所有真实记录至多具有一个数位值。将各包按照数位稳定排序后，数位相同的记录在块内和跨块时均保持原有次序。空包排在最后。图~\ref{fig:sort-packs} 给出两个块进行一趟排序的情形。

\begin{figure}[!hbp] \centering \begin{tabular}{@{}l@{\qquad}c@{\qquad}c@{}}
\toprule
输入块中的原次序 & 数位为 $0$ 的包 & 数位为 $1$ 的包\\
\midrule
$a_1\, b_0\, c_1\, d_1$ & $[b,\cdot]$ & $[a,c]\ [d,\cdot]$\\
$e_0\, f_1\, g_0\, h_0$ & $[e,g]\ [h,\cdot]$ & $[f,\cdot]$\\
\bottomrule
\end{tabular} \caption{一趟基数排序，每包容纳两条记录。下标表示本趟排序依据的数位，圆点表示虚拟槽位。对整包稳定排序，相当于先收集中列，再收集右列，每列都自上而下读取。读出真实记录，得到 $b,e,g,h,a,c,d,f$。} \label{fig:sort-packs} \end{figure}

每一块预留固定的 $O((2^s+1)P)$ 个额外槽位，既足以填充所有真实连续段，也能保留原有虚拟记录，并将总长度向上补齐至 $P$ 的整倍数。扫描各段，求出长度，向上取整，并给所有原有记录（包括原有虚拟记录）分配目的位置。第二次局部带标记置换负责核验这些目的位置。再追加新的虚拟记录，要求其猜测的目的位置严格递增。它们只能进入尚未使用的位置，因此恰有一种合法填法。为凑成二的幂而临时追加的填充，在调用结束后移除。全局稳定排序前扫描每一包以取得其数位。读取真实次序时，忽略全部虚拟记录，包括包内的虚拟记录。

在所有趟次中，列表大小始终为 $O(n)$。在一趟中，其乘法增长因子为 $1+O(2^{-s}P)$，最后一块另有 $O(2^{2s}+2^sP)$ 的加法成本。在 $O(h/s)$ 趟中，这些乘法因子的积有界，因为 $P$ 关于 $h$ 为多项式有界，且 $s=\lfloor h^{1/a}\rfloor$。由于 $s=o(h)$，累积的加法成本为 $o(2^h)$。

每次局部调用的地址长度为 $2s+O(1)$。把记录宽度补到 $h$ 的一个足够大的常数倍，即可满足第一个宽度阈值，因此每个槽位的成本为 $O(hs^{1/2})$。全局排序有 $O(n/P)$ 个包，每包宽 $O(hP)$ 比特，地址长度为 $h+O(\log h)$。因 $hP\ge h^B$，按常数倍填充即可满足第二个阈值，全局成本为 $O(nh^{1+1/(2a)})$。扫描、复制和次序检查也满足同一上界。所有趟次总共花费 \[
            O\bigl((h/s)\,nh^{1+1/(2a)}\bigr)
               =O\bigl(nh^{\,2-1/(2a)}\bigr).
\]

为实现 $\mathrm{Tag}(h,h)$，将 $n$ 条源记录与 $n$ 条候选目的记录合并。赋予每条源记录其声称的目的地址，给每条目的记录指定其规定位置，然后将整个列表按目的地址作基数排序，地址相同时源记录在前、目的记录在后。跳过虚拟记录，扫描真实记录。按顺序对每个目的地址要求先有一条源记录，紧接着有一条有效载荷相同的目的记录。一个 $O(h)$ 比特寄存器即已足够。该扫描拒绝源标记重复或缺失的情形，也无须重新排列最后的虚拟槽位。

因此 $\mathrm{Tag}(h,h)$ 的指数为 $p=2-1/(2a)$，且 $1<p<2$。对宽度 $ch$，将有效载荷分到有界多次使用相同标记的带标记置换调用中，必要时予以填充。利用从稳定排序到带标记置换的归约即得通常宽度情形的结论，从而完成定理~\ref{thm:sort} 的证明。
